\section{First Graphics Bring-Up: A Null Pointer With a Documented Cause}
\label{sec:graphics1}

The installation is complete and the executable is the publisher's own. What
follows is the second half of the project, which took considerably longer than
the first, and which is entirely about the graphics stack.

\subsection{Establishing that the reconstruction is actually the game}

Before diagnosing rendering, it was worth confirming that the 54\,GB tree is a
coherent game rather than a coherent-looking pile. Four observations settle it,
and all four come from the game's own startup behaviour rather than from our
inspection:

\begin{itemize}
\item The engine locates and reads its encrypted package definition,
      \verb|packagedefinition.txt| at 35{,}284\,bytes, which is the file the
      Glacier engine consults to find every other asset. A wrong or truncated
      copy stops startup immediately.
\item The crash-telemetry component initialises and writes
      \verb|crash_metrics.dat|, 25{,}240\,bytes.
\item NVIDIA's Streamline framework loads seven \verb|sl.*| libraries and
      \emph{verifies its own digital signatures on each of them}, which is a
      third-party integrity check passing on our reconstruction, performed by
      code we did not write and cannot influence.
\item Streamline then allocates its device-side and host-side pools,
      \verb|sl.chi.heapGPU| and \verb|sl.chi.heapCPU|.
\end{itemize}

The game reaches its rendering path. Everything after this point is a graphics
problem, not a packaging problem.

\subsection{The bundled translation layer is too old}

Direct3D~12 is translated to Vulkan by vkd3d-proton~\cite{vkd3d}. The
system Wine ships a much older vkd3d, and its failure is immediate and legible:

\begin{lstlisting}
Enhanced barriers: 0
\end{lstlisting}

The game requires enhanced barriers, a Direct3D~12 feature for expressing
resource state transitions, and refuses to proceed without them. This is a clean,
early, well-signposted failure---the kind one wants.

Replacing the translation layer with vkd3d-proton 2.14.1 and the Direct3D~11
path with DXVK~\cite{dxvk} changes the log to:

\begin{lstlisting}
DX Ultimate supported!
DXR 1.1 support enabled
Enabling support for SM 6.8
\end{lstlisting}

Ray tracing, mesh shaders, and Shader Model 6.8 all negotiate successfully. The
game then crashes.

\subsection{The crash: a null vtable, disassembled}

The fault is a read from address zero, and the three instructions around it say
exactly what happened:

\begin{lstlisting}
140cb497c:  mov  0x8(%rax),%rcx   ; load an interface ptr
140cb4997:  mov  (%rcx),%rax      ; load its vtable  <-- FAULT
140cb49a8:  call *(%rax)          ; call through it
\end{lstlisting}

This is the canonical shape of an unchecked COM interface query. The game asked
for an interface, did not check the return value, stored the null it received,
and called a method on it one basic block later.

\subsubsection{Which interface}

The requested interface is \verb|ID3D12GraphicsCommandList10|, IID
\verb|{7013c015-d161-4b63-a08c-238552dd8acc}|. It is a recent addition to the
Direct3D~12 interface family. vkd3d-proton 2.14.1 does not implement it and
correctly returns \verb|E_NOINTERFACE|; the game ignores that and dereferences
the null.

\subsubsection{Verifying support without running anything}

A useful and very cheap technique: a COM interface's presence in a translation
layer can be checked by searching the binary for the IID in its
\emph{little-endian in-memory byte order}, which is not the order the GUID is
printed in.

\begin{lstlisting}
grep -qaF $'\x15\xc0\x13\x70\x61\xd1\x63\x4b'\
$'\xa0\x8c\x23\x85\x52\xdd\x8a\xcc' d3d12core.dll
\end{lstlisting}

The first three fields byte-swap; the last two do not. Getting this wrong
produces a confident false negative.

A second false-negative source cost real time and is worth stating plainly:
inside a Proton distribution the Direct3D~12 libraries live under
\verb|files/lib/wine/vkd3d-proton/x86_64-windows|. \textbf{That is
\texttt{lib}, not \texttt{lib64}.} A \verb|lib64| path exists and contains other
things, so a search there returns nothing and looks like an answer.

A third: Wine ships its own \verb|d3d12.dll| of roughly 124\,KB, which is a thin
forwarding stub and is \emph{not} vkd3d-proton. Confirming the size distinguishes
them at a glance.

\begin{table}[htbp]
\caption{Installed translation-layer libraries, for size comparison}
\label{tab:dlls}
\centering
\footnotesize
\begin{tabular}{@{}lr@{}}
\toprule
\textbf{Library} & \textbf{Bytes} \\
\midrule
\verb|d3d12.dll| & 159{,}758 \\
\verb|d3d12core.dll| & 5{,}185{,}550 \\
\verb|d3d11.dll| & 6{,}209{,}550 \\
\verb|d3d10core.dll| & 290{,}830 \\
\verb|dxgi.dll| & 4{,}435{,}982 \\
\verb|nvapi64.dll| & 1{,}974{,}286 \\
\verb|nvofapi64.dll| & 1{,}384{,}462 \\
\bottomrule
\end{tabular}
\end{table}

\subsection{An unrelated fault in the same session: divide by zero in the display path}

Independently of the interface problem, the game faulted while enumerating
display modes. The log sequence is self-explanatory once the hardware topology
from Section~\ref{sec:system} is in mind:

\begin{lstlisting}
Found monitors not associated with any adapter,
  using fallback
readMonitorEdidFromKey: Failed to get EDID
  reg key size
\end{lstlisting}

The discrete GPU has no attached display. Enumerating modes for the rendering
adapter therefore yields an empty set, and the engine divides by a refresh rate
taken from that empty set.

The fix is to give the application a display to find: a Wine virtual desktop, at
which point mode enumeration returns a single well-formed entry and the division
is well defined. This costs nothing and is required on any Optimus-class laptop.

\subsection{The evidence that reversed an hour of reasoning}

The log reported the graphics adapter as vendor \verb|0x1002|, device
\verb|0x73df|---an AMD Radeon. On a machine with two GPUs, running a test
intended to use the NVIDIA card, this reads unambiguously as ``the override did
not take.''

It is the opposite. \textbf{DXVK deliberately reports an AMD device while
running on NVIDIA.} Two options, \verb|dxgi.nvapiHack| and
\verb|dxgi.hideNvidiaGpu|, default to enabled and exist because a number of
shipped games take a worse code path when they detect an NVIDIA adapter through
DXGI. Seeing the AMD identifiers is therefore \emph{positive evidence that the
NVIDIA card is in use}.

This is worth generalising. \textbf{A compatibility layer's job is to lie
convincingly, so its self-reports are the worst available evidence about what
hardware is in use.} Ground truth must come from outside the layer---in this
project, from \verb|nvidia-smi --query-compute-apps|, which reports the process
list the kernel driver actually has.

\subsection{Hypotheses eliminated at this stage}

\begin{table}[htbp]
\caption{Eliminated during first bring-up}
\label{tab:elim-g1}
\centering
\footnotesize
\begin{tabularx}{\columnwidth}{@{}L{2.9cm}X@{}}
\toprule
\textbf{Hypothesis} & \textbf{How it died} \\
\midrule
Vendor-gated crash & Spoofing the vendor as Intel (\verb|8086|) and as NVIDIA
                     (\verb|10de|) both crash at the \emph{identical} address \\
AMD helper library at fault & \verb|amd_ags_x64.dll| is a hard static import;
                              removing it yields \verb|c0000135| before any
                              rendering \\
Never reached rendering & The log shows the swapchain created and presented \\
\bottomrule
\end{tabularx}
\end{table}

The identical crash address under three different spoofed vendors is the
strongest of these: a vendor-conditional code path would branch, and branching
code does not fault at one address.

\subsection{Where this leaves the investigation}

Upgrading the translation layer resolves the missing interface, and the virtual
desktop resolves the display-mode fault. The game then stops crashing---and
starts doing something much harder to diagnose, which is the subject of
Section~\ref{sec:deadlock}. Before that, Section~\ref{sec:egl} deals with a
finding that had to be retracted, because for three sessions the discrete GPU was
believed to be incapable of Vulkan at all.
